\scanpage{12}
conclusion.

\uline{\deleted{\unsure{N}}Remarque} Soit $A$ un schéma abélien sur $V$,
$B'$ un sous-schéma abélien de $A'=A_K$,
Il résulte de ce qui précède que $B'$
est de la forme $B_K$, où $B$ est un
schéma abélien sur $V$, et où l'immersion
$B'\mathrel{-}A'$ est \unsure{induit|induite} par un morphisme
$B\longrightarrow A$ (ce dernier fait résulte
aussi du th. de Weil). \deleted{Mais} Soit
$H$ le sous-schéma fermé de $A$ image de $B$,
on voit \deleted{\unsure{facilement}}\inserted{tout de suite} que $H=\overline{B'}$ (car $H$
induit $B'$, \deleted{et} donc $H$ majore $\overline{B'}$,
or \deleted{\unsure{$B'$|$B$}} \unsure{l'|$1^\circ$}image inverse de $\overline{B'}$ est un
sous-schéma fermé de $B$ \unsure{donnant|dominant} $B\otimes K$, donc
identique à $B$, donc $\overline{B'}$ \deleted{est} majore $H$ :).
C'est donc un sous-schéma en groupes
fermé de $A$, \inserted{plat sur $V$.} (prop. 1, cor 4). C'est
un schéma abélien sss $B\longrightarrow B'$ est
un isom. (grâce au th de Weil p.ex),
i.e. $B\longrightarrow A$ est un monomorphisme,
i.e. son \uline{noyau} est \deleted{\illegible{1}} trivial. --
D'après \unsure{Shimura} - Koizumi, il n'en est
% archive sheet number: 11

